\documentclass[journal]{IEEEtran}

% =========================================================================
% PACKAGES & CONFIGURATION
% =========================================================================
\usepackage{amsmath,amssymb,amsfonts,bm}
\usepackage{graphicx}
\usepackage{booktabs}
\usepackage{multirow}
\usepackage{array}
\usepackage{makecell}
\usepackage{cite}
\usepackage{url}
\usepackage{microtype}
\usepackage{balance}
\usepackage{xcolor}
\usepackage[pagebackref=false,breaklinks=true,colorlinks=true,citecolor=blue,linkcolor=blue,urlcolor=blue]{hyperref}

% Correct bad hyphenation here
\hyphenation{op-ti-cal net-works semi-conduc-tor dis-en-tan-gle-ment pa-thol-o-gy}

\begin{document}

% =========================================================================
% TITLE & AUTHOR BLOCK
% =========================================================================
\title{DFU-WAGNER net: Multi-Task Diabetic Foot Ulcer Characterization via Hybrid Feature Fusion, Subspace Disentanglement, and Uncertainty-Gated Domain Alignment}

\author{
    Priyanka~Saxena,~
    Aniketh~Vuruputuri,~
    Nikunj~Jayesh,~
    U.~V.~S.~Pavan~Kumar,~
    Aditi,~
    and~Jagadhabhiram~Medishetty%
    \thanks{P. Saxena, A. Vuruputuri, N. Jayesh, U. V. S. Pavan Kumar, Aditi, and J. Medishetty are with the Department of Computer Science and Engineering / Department of Biomedical Engineering (e-mail: corresponding.author@research.org).}%
}

% Paper headers
\markboth{IEEE Transactions on Medical Imaging,~Vol.~XX, No.~X, October~2026}%
{Saxena \MakeLowercase{\textit{et al.}}: DFU-WAGNER net: Multi-Task Diabetic Foot Ulcer Characterization}

\maketitle

% =========================================================================
% ABSTRACT
% =========================================================================
\begin{abstract}
Diabetic foot ulcers (DFUs) represent a severe global healthcare burden requiring concurrent assessment along two clinically distinct diagnostic axes: microbiological/etiological pathology (\textit{None}, \textit{Infection}, \textit{Ischaemia}, \textit{Both}) and anatomical tissue invasion depth (Wagner Grades 0--4). However, automated diagnostic algorithms face critical clinical and computational bottlenecks: existing public cohorts are fragmented into isolated annotation regimes; severe class imbalance suppresses limb-threatening ischaemia detection (3.81\% prevalence in DFUC2021); and cross-cohort acquisition differences induce negative transfer across disparate hospital sources. Furthermore, cryptographic audit of public repositories reveals uncurated cross-dataset duplicate overlap resulting from web-origin scraping, requiring rigorous split quarantine.

To overcome these challenges, this paper presents \textbf{DFU-WAGNER net}, a unified multi-task deep representation learning framework. DFU-WAGNER net couples an \textbf{EfficientNetV2-S} convolutional stream with a \textbf{Swin-Transformer Tiny} stream through symmetrical 1,280-dimensional linear projections, capturing localized epidermal textures alongside macroscopic spatial foot topology within a fused 2,560-dimensional representation. Rather than utilizing monolithic feature sharing, a dedicated three-way \textbf{TaskDisentangler} partitions the latent space into an $L_2$-normalized shared anatomical subspace ($z_{\text{shared}} \in \mathbb{R}^{768}$) and two task-isolated private subspaces ($z_{\text{path\_private}} \in \mathbb{R}^{512}$ and $z_{\text{depth\_private}} \in \mathbb{R}^{512}$), regularized by a pairwise cosine \textbf{Orthogonality Loss} ($\mathcal{L}_{\text{orth}}$). 

To prevent cross-cohort distribution collapse, we introduce \textbf{Uncertainty-Gated Maximum Mean Discrepancy (UG-MMD)}, which aligns \textit{only} the shared subspace between cohorts using multi-scale Gaussian kernels dynamically modulated by detached predictive entropy confidence weights. Long-tailed pathology classification is optimized via \textbf{Class-Balanced Focal Loss} and \textbf{Hard-Pair Tangent LMCoT Metric Learning} on $z_{\text{shared}}$, while anatomical ulcer depth is modeled using \textbf{Consistent Rank Logits (CORAL)} ordinal cumulative cutoffs. 

The primary 59,783,374-parameter architecture was evaluated across the DFUC2021 benchmark ($n = 5,955$ labeled samples) and a deduplicated Wagner cohort ($n = 11,491$ training samples). On the 20\% development validation split ($n = 1,191$), our primary 45-epoch model achieves a Macro-F1 of \textbf{0.9201} and an overall accuracy of \textbf{90.18\%}, elevating \textit{Ischaemia} F1-score to \textbf{0.9149} with \textbf{93.48\%} sensitivity. Concurrently, the network achieves an ordinal Quadratic Weighted Kappa (QWK) of \textbf{0.7765}, adjacent depth accuracy of \textbf{91.42\%}, and Mean Absolute Error (MAE) of \textbf{0.4290} on the Wagner grading validation set ($n = 303$). A $5 \times 4$ exploratory joint probability distribution models concurrent risk under conditional independence. Systematic ablations validate that active hard-pair mining provides a $+4.31\%$ Macro-F1 gain, while multi-task trade-offs and cross-dataset duplicate integrity are formally quantified.
\end{abstract}

% =========================================================================
% INDEX TERMS
% =========================================================================
\begin{IEEEkeywords}
Diabetic Foot Ulcer, Multi-Task Learning, Subspace Disentanglement, Ordinal Regression, Vision Transformer, Metric Learning, Domain Adaptation, Explainable AI.
\end{IEEEkeywords}

\IEEEpeerreviewmaketitle

% =========================================================================
% SECTION I: INTRODUCTION
% =========================================================================
\section{Introduction}
\IEEEPARstart{D}{iabetic} Foot Ulcers (DFUs) are among the most debilitating, costly, and life-threatening chronic complications of diabetes mellitus, affecting approximately 18.6 million individuals globally each year \cite{yap2022development, schaper2020practical}. Peripheral sensory neuropathy impairs protective pain feedback, causing repetitive biomechanical stress that breaks down fragile plantar skin \cite{schaper2020practical}. Over 50\% of ulcers become clinically infected, and approximately 20\% of moderate-to-severe infections lead to lower-extremity amputation \cite{schaper2020practical}. Epidemiological studies indicate that the five-year post-amputation mortality rate exceeds 50\%, a prognosis exceeding several aggressive malignancies \cite{yap2022development, armstrong1998validation}.

\subsection{Clinical Motivation}
In clinical wound practice, patient management is governed by two complementary diagnostic axes:
\begin{enumerate}
    \item \textbf{Microbiological and Etiological Pathology}: Evaluating microbial infection, arterial perfusion insufficiency (ischaemia resulting from peripheral arterial disease), or their concurrent presentation (neuro-ischaemic ulceration) \cite{goyal2020recognition, armstrong1998validation}. This axis determines urgent medical therapy, including systemic antimicrobial regimens and vascular revascularization.
    \item \textbf{Anatomical Invasion Depth and Necrosis}: Assessing the structural penetration depth through epidermis, dermis, tendon, joint capsule, periosteum, and bone (osteomyelitis) or gangrene, traditionally staged via the Meggitt-Wagner grading system (Grades 0--4) \cite{wagner1981dysvascular}. This axis governs surgical intervention, debridement, and compartment drainage.
\end{enumerate}

In clinical prognostication, systems such as the University of Texas (UT) Diabetic Wound Classification System \cite{armstrong1998validation} conceptually pair depth stages with infection/ischaemia categories. An ulcer presenting with superficial skin loss but critical ischaemia requires urgent revascularization, whereas an ulcer with deep joint capsule exposure but clean perfusion requires aggressive debridement. Evaluating either dimension in isolation provides an incomplete diagnostic picture \cite{armstrong1998validation}.

\subsection{Computer Vision Motivation and Public Cohort Fragmentation}
Digital photography acquired via smartphones and point-of-care cameras provides a non-invasive medium for wound monitoring and telemedicine \cite{goyal2019robust, yap2022development}. Over the past decade, Convolutional Neural Networks (CNNs) \cite{goyal2020dfunet, goyal2019robust, goyal2020recognition} and Vision Transformers (ViTs) \cite{galdran2021convolutional, karthik2025attention} have advanced automated image analysis.

However, translation to unified multi-task assessment is hindered by \textbf{clinical taxonomy fragmentation} in public benchmarks \cite{yap2022development}:
\begin{itemize}
    \item The \textbf{DFUC2021 Benchmark} \cite{cassidy2021diabetic} provides 5,955 labeled wound patches annotated exclusively for infection and ischaemia (\textit{None}, \textit{Infection}, \textit{Ischaemia}, \textit{Both}), with zero annotations regarding ulcer depth or osteomyelitis.
    \item The public \textbf{Wagner Ulcer Cohorts} \cite{rostami2024deep, zhou2025deep} provide photographs annotated exclusively according to the 5-point Meggitt-Wagner scale (Grades 0 to 4), without microbiological or vascular labels.
\end{itemize}

No public clinical benchmark provides paired, co-annotated labels for both taxonomies. Previous computer vision models have therefore functioned as isolated single-task classifiers \cite{toofanee2023dfu, galdran2021convolutional, rostami2024deep, zhou2025deep}.

\subsection{Literature-Supported Methodological Challenges}
A rigorous review of the existing DFU deep learning literature reveals four primary bottlenecks:
\begin{enumerate}
    \item \textbf{Class Imbalance in Long-Tailed Pathology}: In DFUC2021, the critical minority class---\textit{Ischaemia}---accounts for only 3.81\% ($n = 227$) of labeled training samples, yielding an 11.26:1 majority-to-minority ratio \cite{cassidy2021diabetic}. Standard cross-entropy optimization causes gradient starvation, leading to high false-negative rates on ischaemic limbs \cite{toofanee2023dfu, galdran2021convolutional}.
    \item \textbf{Disregard for Ordinal Biophysical Geometry}: Previous deep learning studies addressing Wagner grading \cite{rostami2024deep, zhou2025deep} formulate the task as nominal multi-class classification, penalizing a minor adjacent error (Grade 1 vs. Grade 2) identically to a severe diagnostic failure (Grade 0 vs. Grade 4). Because tissue necrosis follows an ordinal continuum, nominal models fail to bound Mean Absolute Error (MAE) or optimize Quadratic Weighted Kappa (QWK) \cite{cao2020rank, cohen1968weighted}.
    \item \textbf{Monolithic Feature Interference and Negative Transfer}: When multi-task networks share representations across disparate tasks, monolithic parameter sharing forces conflicting gradients through a single feature vector \cite{caruana1997multitask, bousmalis2016domain}. Wagner depth gradients can corrupt representations required for superficial erythema and vascular discoloration.
    \item \textbf{Cross-Dataset Acquisition Disparities and Scraping Overlap}: DFUC2021 images consist of tightly cropped wound patches from UK NHS hospital trusts \cite{cassidy2021diabetic}, whereas Wagner images comprise diverse clinical photography from international cohorts \cite{rostami2024deep}. Passing both through unaligned backbones induces domain drift. Furthermore, cryptographic audit reveals that public web scraping has introduced exact duplicate images across datasets, demanding rigorous data quarantine.
\end{enumerate}

\subsection{Proposed Approach: DFU-WAGNER net}
To address these limitations, we propose \textbf{DFU-WAGNER net}, an approximately 60-million parameter multi-task deep representation learning architecture. DFU-WAGNER net integrates an \textbf{EfficientNetV2-S} CNN branch (extracting localized textural and edge features) with a \textbf{Swin-Transformer Tiny} branch (extracting global spatial context) via symmetrical 1,280-dimensional linear projections, forming a fused 2,560-dimensional hybrid representation.

Crucially, DFU-WAGNER net introduces a three-way \textbf{TaskDisentangler} module that decomposes the fused representation into three $L_2$-normalized subspaces:
\begin{itemize}
    \item A shared anatomical subspace: $z_{\text{shared}} \in \mathbb{R}^{768}$.
    \item A pathology-private subspace: $z_{\text{path\_private}} \in \mathbb{R}^{512}$.
    \item A depth-private subspace: $z_{\text{depth\_private}} \in \mathbb{R}^{512}$.
\end{itemize}

An $\mathcal{O}(B \cdot D)$ cosine \textbf{Orthogonality Regularization Loss} ($\mathcal{L}_{\text{orth}}$) enforces linear independence between shared and private channels. Distribution alignment between the DFUC2021 and Wagner cohorts is enforced \textit{exclusively} on $z_{\text{shared}}$ via \textbf{Uncertainty-Gated Maximum Mean Discrepancy (UG-MMD)}, weighted by detached model predictive entropy. Pathology classification is supervised via \textbf{Class-Balanced Focal Loss} and \textbf{Hard-Pair Tangent LMCoT Metric Learning} on $z_{\text{shared}}$, while Wagner depth is regressed via \textbf{Consistent Rank Logits (CORAL)} ordinal binary cutoffs.

\subsection{Summary of Scientific Contributions}
The verified contributions of this work are:
\begin{enumerate}
    \item \textbf{Unified Multi-Task Hybrid Architecture}: An asymmetric dual-backbone network combining EfficientNetV2-S and Swin-Transformer Tiny via symmetrical 1,280-D projections, concurrently predicting 4-class pathology and 5-grade Wagner ordinal depth without requiring co-annotated multi-label training data.
    \item \textbf{Shared-Private Subspace Disentanglement}: A 3-way \texttt{TaskDisentangler} partitioning fused representations into $z_{\text{shared}} \in \mathbb{R}^{768}$, $z_{\text{path\_private}} \in \mathbb{R}^{512}$, and $z_{\text{depth\_private}} \in \mathbb{R}^{512}$, regularized by cosine orthogonality $\mathcal{L}_{\text{orth}}$ to prevent cross-task feature corruption.
    \item \textbf{Hard-Pair Tangent LMCoT Metric Learning}: A monotonic tangent distance $\tan(\theta/2)$ on $z_{\text{shared}}$ with active batch-level hard-pair mining across clinical decision boundaries to counter minority class collapse.
    \item \textbf{Uncertainty-Gated Multi-Scale RKHS-MMD Alignment}: A distribution alignment objective aligning \textit{only} the shared subspace across cohorts, dynamically gated by detached predictive entropy to prevent noisy or out-of-distribution representations from distorting the aligned manifold.
    \item \textbf{Rank-Consistent Ordinal Depth Modeling}: Integration of CORAL ordinal cumulative cutoffs for Wagner grading to preserve biological continuum and directly optimize QWK and MAE.
    \item \textbf{Cryptographic Cross-Dataset Overlap Audit}: Programmatic disclosure and quarantine of 728 cross-dataset duplicate matches resulting from public web scraping, providing open manifests for reproducible, leak-free multi-task benchmarking.
\end{enumerate}

% =========================================================================
% SECTION II: RELATED WORK
% =========================================================================
\section{Related Work}

\subsection{Deep Learning for DFU Classification}
Automated computer-aided diagnosis of diabetic foot ulcers initially focused on patch detection and binary classification (wound vs. healthy skin). Goyal \textit{et al.} \cite{goyal2020dfunet} designed DFUNet, which evaluated convolutional layers with progressive feature reuse to separate ulcer margins from background skin. Subsequent works incorporated region-based object detectors, such as Faster R-CNN and YOLO, to localize ulcers within unconstrained photography \cite{goyal2019robust}. Goyal \textit{et al.} \cite{goyal2020recognition} subsequently compiled a multi-class infection and ischaemia dataset, demonstrating that CNNs could identify surface purulence and arterial discoloration.

\subsection{The DFUC2021 Benchmark and Existing Solutions}
The Diabetic Foot Ulcers Grand Challenge 2021 (DFUC2021) formalized the pathology task into four distinct categories: \textit{None}, \textit{Infection}, \textit{Ischaemia}, and \textit{Both} \cite{cassidy2021diabetic}. The challenge-winning approach by Galdran \textit{et al.} \cite{galdran2021convolutional} evaluated CNNs (ResNeXt50, EfficientNet-B3) and Vision Transformers (DeiT, ViT), observing that standard Transformers trained from scratch struggled with limited sample sizes. Incorporating Sharpness-Aware Minimization (SAM) yielded an official competition hidden-test-set Macro-F1 of 0.6307 \cite{galdran2021convolutional}. Other challenge participants integrated synthetic augmentation and ensemble averaging \cite{cassidy2021diabetic}. However, challenge solutions experienced marked drops on \textit{Ischaemia}, where test F1-scores remained below 0.42.

\subsection{Automated Wagner Grading and Severity Assessment}
The Meggitt-Wagner classification system \cite{wagner1981dysvascular} remains widely utilized for lower-extremity wound penetration depth. Rostami \textit{et al.} \cite{rostami2024deep} evaluated DenseNet backbones across 2,480 images, reporting nominal classification accuracies above 88\%. Zhou \textit{et al.} \cite{zhou2025deep} coupled Mask2Former with a Swin Transformer to segment boundaries and predict Wagner grades 1 through 4. However, existing Wagner studies universally formulate depth grading as nominal multi-class classification \cite{rostami2024deep, zhou2025deep}, discarding the ordinal structure of tissue breakdown and penalizing adjacent grade boundaries identically to multi-grade errors.

\subsection{Hybrid CNN-Transformer Medical Image Classification}
Recent medical imaging literature highlights complementary strengths of CNNs and Vision Transformers \cite{tan2021efficientnetv2, liu2021swin}. Convolutional networks, particularly EfficientNetV2 \cite{tan2021efficientnetv2}, provide parameter-efficient extraction of localized translational features, including micro-vascular erythema, wound exudate, and granular tissue margins. Conversely, Vision Transformers---notably the Swin Transformer \cite{liu2021swin}---utilize shifted local attention windows to model long-range spatial context, such as anatomical foot positioning and bilateral symmetry, with linear computational scaling $\mathcal{O}(N)$. Hybrid architectures combining both streams have demonstrated superior robustness across clinical classification benchmarks \cite{karthik2025attention, mehta2022mobilevit}.

\subsection{Siamese Representation and Metric Learning}
Metric learning optimizes the geometry of the latent embedding space directly by enforcing that semantically identical images cluster tightly together while dissimilar images are separated by a margin \cite{bromley1993signature, chopra2005learning}. The foundational \textbf{DFU-SIAM} framework by Toofanee \textit{et al.} \cite{toofanee2023dfu} introduced Siamese learning to DFU classification by coupling EfficientNetV2-S with BEiT under a Large Margin Cotangent (LMCoT) loss, reporting a Macro-F1 of 0.6228 on DFUC2021. However, the original formulation was restricted to single-task pathology, lacked active hard-pair mining, and exhibited cotangent polarity sensitivity near boundaries.

\subsection{Multi-Task Learning and Representation Disentanglement}
Multi-task learning (MTL) leverages shared representations across related tasks to improve data efficiency and generalization \cite{caruana1997multitask}. However, naive multi-task networks frequently suffer from negative gradient interference when joint tasks require conflicting feature representations. Bousmalis \textit{et al.} \cite{bousmalis2016domain} introduced Domain Separation Networks (DSNs), which partition representations into shared and private subspaces regularized by an orthogonality penalty. Disentangled representations ensure that task-specific features (such as osteomyelitis bone invasion) do not degrade representations optimized for independent tasks (such as surface erythema) \cite{bousmalis2016domain}.

\subsection{Domain Adaptation and Distribution Alignment}
Transferring representations across clinical cohorts is hampered by domain shift---differences in camera hardware, lighting, patient skin tone distributions, and institutional imaging protocols \cite{yap2022development}. Maximum Mean Discrepancy (MMD) \cite{gretton2012kernel, long2015learning} maps representations into a Reproducing Kernel Hilbert Space (RKHS) and minimizes the statistical distance between empirical sample distributions. Kendall \textit{et al.} \cite{kendall2018multi} demonstrated that incorporating task uncertainty into loss formulations prevents noisy, out-of-distribution samples from dominating multi-objective optimization.

\subsection{Explainable AI in Medical Imaging}
Given the safety-critical nature of clinical diagnosis, interpretability is essential for clinical adoption. Gradient-weighted Class Activation Mapping (Grad-CAM) \cite{selvaraju2017grad} calculates gradients of target class logits with respect to convolutional feature maps to produce 2D attribution heatmaps. Chefer \textit{et al.} \cite{chefer2021transformer} extended gradient-based attribution to Vision Transformers. In this work, we emphasize that \textbf{attribution heatmaps represent post-hoc gradient saliency and must not be misinterpreted as segmentation masks or definitive causal clinical proofs}.

% =========================================================================
% SECTION III: DATASET & PROBLEM FORMULATION
% =========================================================================
\section{Dataset Integrity and Problem Formulation}

\subsection{DFUC2021 Pathology Dataset}
The pathology classification task utilizes the benchmark dataset established by the \textbf{Diabetic Foot Ulcer Grand Challenge 2021 (DFUC2021)} \cite{cassidy2021diabetic}. The cohort consists of photographic ulcer patches acquired across three National Health Service (NHS) hospital trusts in the United Kingdom.

\subsubsection{Data Integrity and Distribution Audit}
The raw CSV manifest (\texttt{train.csv}) contains 9,949 entries. A programmatic audit revealed:
\begin{itemize}
    \item \textbf{Unlabeled Rows}: Exactly 3,994 entries (40.14\%) contain all-NaN values across all four class columns. These unannotated images were intended by challenge organizers for semi-supervised exploration. In our pipeline, these rows are filtered out prior to supervised training.
    \item \textbf{Labeled Cohort}: Exactly 5,955 images contain verified ground-truth labels.
    \item \textbf{Class Breakdown}:
    \begin{itemize}
        \item \textit{None} (Control): 2,552 images (42.85\%)
        \item \textit{Infection}: 2,555 images (42.91\%)
        \item \textit{Ischaemia}: 227 images (\textbf{3.81\%})---\textit{Severe minority class}
        \item \textit{Both} (Infection + Ischaemia): 621 images (10.43\%)
    \end{itemize}
    \item \textbf{Imbalance Ratio}: The majority-to-minority ratio (\textit{Infection} vs. \textit{Ischaemia}) is \textbf{11.26 : 1.00}.
\end{itemize}

\subsubsection{Development Train / Validation Partitioning}
To evaluate model convergence prior to test evaluation, we implement a stratified 80/20 train/validation split using fixed random seed 42:
\begin{itemize}
    \item \textbf{Training Set} ($n = 4,764$): \textit{None}: 2,042 $|$ \textit{Infection}: 2,044 $|$ \textit{Ischaemia}: 181 $|$ \textit{Both}: 497.
    \item \textbf{Validation Set} ($n = 1,191$): \textit{None}: 510 $|$ \textit{Infection}: 511 $|$ \textit{Ischaemia}: 46 $|$ \textit{Both}: 124.
    \item \textbf{Internal Integrity}: Cryptographic hash audit confirms zero duplicate overlap between internal training and validation partitions.
\end{itemize}

\subsection{Wagner Depth Dataset}
The ulcer depth grading task utilizes the \textbf{Wagner Ulcer Dataset} \cite{rostami2024deep, zhou2025deep}, annotating diabetic foot ulcers across five ordinal stages according to the Meggitt-Wagner grading system \cite{wagner1981dysvascular}:
\begin{itemize}
    \item \textbf{Grade 0}: Pre-ulcerative lesion or healed intact skin.
    \item \textbf{Grade 1}: Superficial ulcer of the epidermis and dermis.
    \item \textbf{Grade 2}: Deep ulcer penetrating to tendon, joint capsule, or bone.
    \item \textbf{Grade 3}: Deep ulcer complicated by abscess, osteomyelitis, or joint sepsis.
    \item \textbf{Grade 4}: Localized gangrene (forefoot or heel necrosis).
\end{itemize}

\subsubsection{Internal Deduplication Audit}
A cryptographic MD5 hash audit across the raw 12,099 images revealed an internal leakage risk:
\begin{itemize}
    \item Exactly 37 base images in \texttt{WagnerDataset/train/Grade 0} were byte-for-byte identical to evaluation images (21 in \texttt{valid/} and 16 in \texttt{test/}).
    \item An additional 111 augmented training images were generated directly from those 37 base images.
    \item To prevent artificial validation inflation without altering raw disk archives, we implemented a manifest-driven exclusion loader (\texttt{wagner\_dedup\_manifest.json}) excluding these 148 images from the training set.
\end{itemize}

\subsubsection{Clean Wagner Cohort Statistics}
Following internal deduplication, the cohort distribution is:
\begin{itemize}
    \item \textbf{Deduplicated Training Set} ($n = 11,491$): Grade 0: 1,852 $|$ Grade 1: 2,240 $|$ Grade 2: 2,345 $|$ Grade 3: 2,744 $|$ Grade 4: 2,310.
    \item \textbf{Preserved Validation Set} ($n = 303$): Grade 0: 21 $|$ Grade 1: 86 $|$ Grade 2: 78 $|$ Grade 3: 40 $|$ Grade 4: 78.
    \item \textbf{Preserved Test Set} ($n = 157$): Grade 0: 16 $|$ Grade 1: 44 $|$ Grade 2: 37 $|$ Grade 3: 18 $|$ Grade 4: 42.
\end{itemize}

\subsection{Cross-Dataset Scraping Overlap Audit}
An exhaustive cross-dataset hash audit revealed that DFUC2021 and Wagner are \textbf{not strictly disjoint cohorts}: exactly \textbf{728 exact MD5 duplicate matches} exist between them due to uncurated web scraping in public releases. 

Cross-referencing the internal split assignments across both cohorts revealed the following overlap distribution:
\begin{itemize}
    \item \textbf{Both in Training Splits}: 435 image pairs.
    \item \textbf{DFUC Validation in Wagner Clean Training}: 119 image pairs.
    \item \textbf{Wagner Validation in DFUC Training}: 92 image pairs.
    \item \textbf{Wagner Test in DFUC Training}: 46 image pairs.
    \item \textbf{Both in Validation Splits}: 21 image pairs.
    \item \textbf{DFUC Validation in Wagner Test}: 15 image pairs.
\end{itemize}

\textbf{Quarantine Protocol}: When training joint multi-task models, images present in one dataset's training split that match the other dataset's validation or test split represent potential cross-task visual leakage. To ensure evaluation integrity, we generated an open reproducibility manifest (\texttt{dataset\_cross\_overlap\_manifest.json} and \texttt{.csv}) and established a leak-free quarantine protocol (\texttt{clean\_leak\_free\_split\_manifest.json}). Under this protocol, 138 DFU training images and 120 Wagner training images are quarantined, resulting in clean, leak-free training partitions of \textbf{4,626 DFU images} and \textbf{11,371 Wagner images}, while preserving the 1,191 DFU validation images, 303 Wagner validation images, and 157 Wagner test images completely untouched and leak-free. Preliminary multi-task baselines reported herein reflect the unquarantined legacy split, while subsequent protocols adopt this quarantine.

\subsection{Patient-Level Independence Limitation}
Neither DFUC2021 nor the Wagner dataset provides patient identification metadata, clinical case numbers, or longitudinal timestamps. While our cryptographic hash deduplication guarantees \textbf{complete image-level and augmentation-level independence} across splits, patient-level independence cannot be mathematically proven. All reported evaluation metrics reflect sample-level generalization across independent photographic instances, consistent with standard practice on these public benchmarks.

% =========================================================================
% SECTION IV: METHODOLOGY
% =========================================================================
\section{Proposed DFU-WAGNER net Methodology}

\subsection{Overall Architectural Pipeline}
Given an input RGB image $x \in \mathbb{R}^{B \times 3 \times 224 \times 224}$, the forward processing pipeline proceeds through five modular stages:
\begin{enumerate}
    \item \textbf{Asymmetric Dual-Backbone Feature Extraction}: Processing $x$ simultaneously through an EfficientNetV2-S CNN branch and a Swin-Transformer Tiny branch.
    \item \textbf{Symmetrical Projection}: Mapping disparate CNN (1,280-D) and Transformer (768-D) feature dimensions into symmetrical 1,280-D vector spaces.
    \item \textbf{Hybrid Feature Fusion}: Concatenating projection outputs into a unified 2,560-D latent vector $\mathbf{h}_{\text{fused}}$.
    \item \textbf{TaskDisentangler Subspace Partitioning}: Decomposing fused features into three strictly $L_2$-normalized subspaces: $z_{\text{shared}} \in \mathbb{R}^{768}$, $z_{\text{path\_private}} \in \mathbb{R}^{512}$, and $z_{\text{depth\_private}} \in \mathbb{R}^{512}$.
    \item \textbf{Selective Task Branching}:
    \begin{itemize}
        \item The Pathology Head receives $\mathbf{v}_{\text{path}} = [z_{\text{shared}}; z_{\text{path\_private}}] \in \mathbb{R}^{1280}$.
        \item The Wagner Head receives $\mathbf{v}_{\text{wagner}} = [z_{\text{shared}}; z_{\text{depth\_private}}] \in \mathbb{R}^{1280}$.
        \item Depth features are \textbf{never} passed to the pathology classifier; pathology features are \textbf{never} passed to the Wagner head.
    \end{itemize}
\end{enumerate}

\subsection{CNN and Transformer Feature Extraction}
The convolutional branch employs \textbf{EfficientNetV2-S} \cite{tan2021efficientnetv2} to capture localized epidermal textures, wound exudate, and micro-vascular margins:
\begin{equation}
    \mathbf{f}_{\text{cnn}} = \text{GAP}\left( \text{Backbone}_{\text{CNN}}(x) \right) \in \mathbb{R}^{1280}
\end{equation}
The transformer branch employs \textbf{Swin-Transformer Tiny} \cite{liu2021swin} with shifted local attention windows to capture global foot geometry and spatial context:
\begin{equation}
    \mathbf{f}_{\text{vit}} = \text{LayerNorm}\left( \text{Backbone}_{\text{Swin}}(x) \right) \in \mathbb{R}^{768}
\end{equation}

\subsection{Symmetrical Feature Projection and Fusion}
To eliminate scale and dimensionality disparity between backbones, both representations are projected into 1,280 dimensions via linear transformations with batch normalization and GELU activations, and concatenated into a unified hybrid representation:
\begin{equation}
    \mathbf{h}_{\text{fused}} = [\text{Proj}_{\text{cnn}}(\mathbf{f}_{\text{cnn}}) \,\|\, \text{Proj}_{\text{vit}}(\mathbf{f}_{\text{vit}})] \in \mathbb{R}^{2560}
\end{equation}

\subsection{Task Disentanglement Subspaces}
To prevent feature interference across disparate clinical tasks, the \texttt{TaskDisentangler} module maps $\mathbf{h}_{\text{fused}}$ into three $L_2$-normalized unit hyperspherical subspaces via parallel projection MLPs:
\begin{equation}
    z_{\text{shared}} \in \mathbb{S}^{767}, \quad z_{\text{path\_private}} \in \mathbb{S}^{511}, \quad z_{\text{depth\_private}} \in \mathbb{S}^{511}
\end{equation}
where $z_{\text{shared}}$ captures cross-cohort transferable wound characteristics, while $z_{\text{path\_private}}$ and $z_{\text{depth\_private}}$ isolate task-specific microbiological and anatomical depth features, respectively.

Task features are selectively constructed by combining shared and private channels: $\mathbf{v}_{\text{path}} = [z_{\text{shared}} \,\|\, z_{\text{path\_private}}] \in \mathbb{R}^{1280}$ and $\mathbf{v}_{\text{wagner}} = [z_{\text{shared}} \,\|\, z_{\text{depth\_private}}] \in \mathbb{R}^{1280}$.

\subsection{Loss Formulations and Multi-Task Objective}

\subsubsection{Class-Balanced Focal Loss ($\mathcal{L}_{\text{path}}$)}
To counter the 11.26:1 class imbalance on \textit{Ischaemia}, supervised pathology classification is optimized via Class-Balanced Focal Loss \cite{cui2019class}:
\begin{equation}
    \mathcal{L}_{\text{path}} = -\frac{1}{N} \sum_{i=1}^N \alpha_{y_i} \left(1 - p_{i, y_i}\right)^\gamma \log\left(p_{i, y_i}\right)
\end{equation}
where class weights $\alpha_c = \frac{1 - \beta}{1 - \beta^{n_c}} / \sum_{k} \frac{1 - \beta}{1 - \beta^{n_k}}$ with $\beta = 0.999$ and focusing parameter $\gamma = 2.0$.

\subsubsection{Consistent Rank Logits ($\mathcal{L}_{\text{wagner}}$)}
Wagner depth grading is decomposed into $K-1=4$ binary cumulative tasks $b_{i, k} = \mathbb{I}(y_i^w > k)$ using CORAL \cite{cao2020rank} with shared linear weight vector $\mathbf{w}$ and independent task biases $\{b_k\}$ to enforce rank consistency:
\begin{equation}
    \mathcal{L}_{\text{wagner}} = -\frac{1}{N (K-1)} \sum_{i=1}^N \sum_{k=0}^{K-2} \left[ b_{i, k} \log \sigma(\hat{y}_{i, k}) + (1 - b_{i, k}) \log \left(1 - \sigma(\hat{y}_{i, k})\right) \right]
\end{equation}
where $\hat{y}_{i, k} = \mathbf{w}^T \mathbf{v}_{\text{wagner}, i} + b_k$.

\subsubsection{Hard-Pair Tangent LMCoT Metric Learning ($\mathcal{L}_{\text{lmcot}}$)}
\label{sec:method_metric}
Metric learning on $z_{\text{shared}}$ applies monotonic tangent distance $d_{\tan}(\mathbf{z}_i, \mathbf{z}_j) = \tan\left(\frac{1}{2} \arccos(\mathbf{z}_i^T \mathbf{z}_j)\right)$ with active batch-level mining of hardest positives $j^*$ and hardest negatives $k^*$:
\begin{equation}
    \mathcal{L}_{\text{lmcot}} = \frac{1}{N} \sum_{i=1}^N \left( \left[d_{\tan}(z_i, z_{j^*})\right]^2 + \left[\max\left(0, m - d_{\tan}(z_i, z_{k^*})\right)\right]^2 \right)
\end{equation}
with margin $m = 2.0$ and distance clamp ceiling $d_{\max} = 10.0$.

\subsubsection{Uncertainty-Gated MMD ($\mathcal{L}_{\text{mmd}}$)}
To align feature distributions across cohorts without corrupting task-specific semantics, multi-scale Gaussian RBF MMD \cite{gretton2012kernel} with bandwidths $\sigma \in \{1, 2, 5, 10\}$ is applied strictly to $z_{\text{shared}}$, weighted by detached predictive entropy confidence $w = \text{clamp}(1 - H(\mathbf{p})/\log C, 0.1, 1.0)$:
\begin{equation}
    \mathcal{L}_{\text{mmd}} = \text{MMD}^2_{\mathbf{w}}\left(Z_p, Z_w\right)
\end{equation}
where $H(\mathbf{p}) = -\sum_{c} p_c \log p_c$ denotes normalized predictive Shannon entropy.

\subsubsection{Subspace Orthogonality Regularization ($\mathcal{L}_{\text{orth}}$)}
To guarantee linear independence between shared and private subspaces within task batches, we penalize mean squared cosine similarities \cite{bousmalis2016domain}:
\begin{equation}
    \mathcal{L}_{\text{orth}} = \frac{1}{2} \left( \frac{1}{N_p} \sum_{i=1}^{N_p} \left( z_{i, s}^T z_{i, p} \right)^2 + \frac{1}{N_w} \sum_{j=1}^{N_w} \left( z_{j, s}^T z_{j, d} \right)^2 \right)
\end{equation}
Notice that cross-private correlation $(z_p^T z_d)^2$ is not evaluated because input samples are processed in task-specific mini-batches where only the active task's private subspace is instantiated.

\subsubsection{Total Combined Objective ($\mathcal{L}_{\text{total}}$)}
The complete multi-task loss is:
\begin{equation}
    \mathcal{L}_{\text{total}} = \lambda_{\text{path}} \mathcal{L}_{\text{path}} + \lambda_{\text{coral}} \mathcal{L}_{\text{wagner}} + \lambda_{\text{lmcot}} \mathcal{L}_{\text{lmcot}} + \lambda_{\text{mmd}} \mathcal{L}_{\text{mmd}} + \lambda_{\text{orth}} \mathcal{L}_{\text{orth}}
\end{equation}
with verified experimental hyperparameter weights $\boldsymbol{\lambda} = [1.0, 2.5, 0.5, 0.1, 0.1]$.

\subsection{Exploratory Joint Bivariate Probability Distribution}
Under the multi-task formulation, DFU-WAGNER net outputs calibrated marginal probabilities for both clinical axes: $\mathbf{p}_{\text{pathology}} \in \mathbb{R}^4$ and $\mathbf{p}_{\text{depth}} \in \mathbb{R}^5$ (derived from CORAL cumulative probabilities $P(W > k)$ via ordinal differencing). 

An exploratory $5 \times 4$ joint probability distribution over the 20 discrete (Wagner $\times$ Pathology) combinations can be computed via outer product under a conditional independence assumption:
\begin{equation}
    P(W = w, P = p \mid x) \approx p_{w}^{\text{depth}} \cdot p_{p}^{\text{pathology}}
\end{equation}
We emphasize that \textbf{this $5 \times 4$ product represents an exploratory bivariate distribution and must not be confused with or equated to the 16-cell University of Texas (UT) diabetic wound classification} (which comprises 4 depth grades $\times$ 4 clinical stages). Because public datasets lack co-annotated UT ground truth, this joint output serves as an exploratory diagnostic visualization rather than a clinically validated prognostic risk score.

\subsection{Explainability via Multi-Branch Grad-CAM}
Visual attribution maps $M_{\text{CAM}}(x)$ are extracted via multi-branch Grad-CAM \cite{selvaraju2017grad} targeting the final convolutional stage of EfficientNetV2-S and the final normalization block of Swin-Tiny. Attribution heatmaps represent post-hoc gradient saliency maps and do not constitute segmentation masks or causal clinical proofs.

% =========================================================================
% SECTION V: EXPERIMENTAL SETUP
% =========================================================================
\section{Experimental Setup}

\subsection{Hardware and Software Infrastructure}
All training runs and benchmarking evaluations are executed on a dedicated local hardware platform:
\begin{itemize}
    \item \textbf{GPU}: NVIDIA GeForce RTX 3050 Laptop GPU (6,144 MB GDDR6 VRAM, Driver 572.16, CUDA 12.8)
    \item \textbf{Host Environment}: Windows 11, Python 3.11.15
    \item \textbf{Deep Learning Framework}: PyTorch 2.6.0 with CUDA 12.4 runtime (\texttt{2.6.0+cu124})
    \item \textbf{Mixed Precision}: Automatic Mixed Precision (AMP FP16) utilizing \texttt{torch.amp.GradScaler('cuda', init\_scale=1024.0)}
\end{itemize}

\subsection{Training Optimization and Schedules}
\begin{itemize}
    \item \textbf{Optimizer}: AdamW with decoupled weight decay.
    \item \textbf{Differential Parameter Groups}:
    \begin{itemize}
        \item Backbone parameters (EfficientNetV2-S + Swin-Tiny): Learning rate $\eta_{\text{backbone}} = 2 \times 10^{-5}$, weight decay $= 1 \times 10^{-3}$.
        \item Projection, Disentangler, and Task Heads: Learning rate $\eta_{\text{head}} = 4 \times 10^{-4}$, weight decay $= 1 \times 10^{-3}$.
    \end{itemize}
    \item \textbf{Learning Rate Schedule}: 3-epoch linear warmup ($0.1 \times \eta \to 1.0 \times \eta$) followed by Cosine Annealing over 45 epochs with $\eta_{\min} = 1 \times 10^{-7}$.
    \item \textbf{Gradient Norm Clipping}: Maximum gradient norm threshold $= 1.0$ (unscaled prior to clipping).
    \item \textbf{Batch Size}: Batch sizes of 32 (A0-Path, A0-Wag, A1-Full, A2-NoLMCoT) and 16 (A3-NoMMD, A4-NoOrth, A6-StdCE, A7-NoHardMining) under 4 DataLoader workers.
    \item \textbf{Interleaved Training}: Interleaved DataLoader pairing 6,000 class-balanced DFUC pairs with interleaved Wagner batches per epoch.
\end{itemize}

\subsection{Data Augmentation Pipeline}
During training, images are resized to $224 \times 224$ and subjected to a clinically tailored augmentation pipeline:
\begin{itemize}
    \item Random Affine (rotation $\pm 15^\circ$, translation $\pm 10\%$, scaling $[0.9, 1.1]$, shear $\pm 10^\circ$).
    \item Random Horizontal Flip ($p = 0.5$).
    \item Color Jitter (brightness $\pm 0.2$, contrast $\pm 0.2$, saturation $\pm 0.2$, hue clamped to zero to prevent altering erythema tones).
    \item Random Autocontrast Stretching ($p = 0.2$).
    \item Random Erasing / Cutout ($p = 0.25$, scale $[0.02, 0.2]$).
    \item Normalization using ImageNet channel statistics ($\mu = [0.485, 0.456, 0.406]$, $\sigma = [0.229, 0.224, 0.225]$).
\end{itemize}

\subsection{Evaluation Metrics}
\begin{enumerate}
    \item \textbf{DFUC2021 Pathology Classification}:
    \begin{itemize}
        \item \textbf{Macro-F1}: Unweighted harmonic mean of precision and recall across all 4 classes.
        \item \textbf{Per-Class F1-Scores}: Tracking of \textit{Ischaemia} ($n = 46$) and \textit{Both} ($n = 124$) minority classes.
        \item \textbf{Accuracy}, \textbf{Precision}, and \textbf{Unweighted Average Recall (UAR)}.
    \end{itemize}
    \item \textbf{Wagner Depth Ordinal Grading}:
    \begin{itemize}
        \item \textbf{Quadratic Weighted Kappa (QWK)}: Measures ordinal agreement penalizing off-diagonal errors quadratically.
        \item \textbf{Mean Absolute Error (MAE)}: $\frac{1}{N} \sum_{i=1}^N |\hat{y}_i^w - y_i^w|$ in ordinal grade units.
        \item \textbf{Adjacent Accuracy}: Percentage of predictions satisfying $|\hat{y}_i^w - y_i^w| \le 1$.
        \item \textbf{Exact Accuracy}.
    \end{itemize}
\end{enumerate}

% =========================================================================
% SECTION VI: RESULTS
% =========================================================================
\section{Results}

\subsection{Pathology Classification Performance}
Table~\ref{tab:dfuc_comparison} summarizes performance across the DFUC2021 pathology benchmark. To preserve evaluation integrity, the table explicitly delineates published official competition hidden-test-set evaluations from our internal development validation cohort evaluations.

On the 20\% held-out development validation split ($n = 1,191$), our primary 45-epoch DFU-WAGNER net model achieves a Macro-F1 score of \textbf{0.9201} and an overall accuracy of \textbf{90.18\%}. Most notably, DFU-WAGNER net achieves an F1-score of \textbf{0.9149} on the critical minority class (\textit{Ischaemia}) with a sensitivity of \textbf{93.48\%} (43/46 cases detected). In comparison, the 25-epoch screening model (A1) achieved a Macro-F1 of 0.9109.

\begin{table*}[t]
\centering
\caption{Performance Overview on the DFUC2021 Pathology Benchmark: Delineating Published Hidden Test Set Results from Internal Validation Cohort Evaluations}
\label{tab:dfuc_comparison}
\begin{tabular}{lllccccc}
\toprule
\textbf{Model / Framework} & \textbf{Backbone Architecture} & \textbf{Evaluation Split} & \textbf{Macro-F1} & \textbf{Accuracy (\%)} & \textbf{Precision} & \textbf{Recall (UAR)} & \textbf{Ischaemia F1} \\
\midrule
\multicolumn{8}{l}{\textit{Part A: Official DFUC2021 Competition Hidden Test Set ($n \approx 5,734$)}} \\
DFUC2021 Baseline \cite{cassidy2021diabetic} & CNN + ViT Ensemble & Official Test Set & 0.6228 & 83.12 & 0.6410 & 0.6120 & 0.3812 \\
Galdran \textit{et al.} \cite{galdran2021convolutional} & ResNeXt50 + DeiT & Official Test Set & 0.6307 & 84.50 & 0.6580 & 0.6240 & 0.4120 \\
Toofanee \textit{et al.} (DFU-SIAM) \cite{toofanee2023dfu} & EfficientNetV2-S + BEiT & Official Test Set & 0.6228 & 83.45 & 0.6390 & 0.6180 & 0.3950 \\
Karthik \textit{et al.} \cite{karthik2025attention} & Swin + EMA Attention & Official Test Set & 0.6285 & 84.10 & 0.6490 & 0.6210 & 0.4010 \\
\midrule
\multicolumn{8}{l}{\textit{Part B: Internal Stratified Development Validation Split ($n = 1,191$)}} \\
DFU-WAGNER net (A1, 25-Epoch Screening) & EfficientNetV2-S + Swin-Tiny & 20\% Val Split & 0.9109 & 88.33 & 0.9022 & 0.9205 & 0.9348 \\
\textbf{DFU-WAGNER net (45-Epoch Primary)} & \textbf{EfficientNetV2-S + Swin-Tiny} & \textbf{20\% Val Split} & \textbf{0.9201} & \textbf{90.18} & \textbf{0.9134} & \textbf{0.9279} & \textbf{0.9149} \\
\bottomrule
\end{tabular}
\end{table*}

\subsection{Per-Class Performance Breakdown}
Table~\ref{tab:per_class} details per-class precision, recall, and F1-score across all four DFUC2021 categories on the validation cohort ($n = 1,191$). The confusion matrix demonstrates strong sensitivity: out of 46 true ischaemic ulcer instances, 43 were successfully diagnosed (93.48\% sensitivity). Furthermore, the combined \textit{Both} category (concurrent infection and ischaemia) achieved 100.0\% recall (124/124) with an F1-score of \textbf{0.9841}.

\begin{table}[h]
\centering
\caption{Detailed Per-Class Pathology Classification Results (DFUC2021 Validation Split, $n=1,191$)}
\label{tab:per_class}
\begin{tabular}{lcccc}
\toprule
\textbf{Pathology Class} & \textbf{Support ($n$)} & \textbf{Precision} & \textbf{Recall} & \textbf{F1-Score} \\
\midrule
None (Control) & 510 & 0.9104 & 0.8569 & 0.8828 \\
Infection & 511 & 0.8785 & 0.9198 & 0.8987 \\
Ischaemia (Minority) & 46 & 0.8958 & 0.9348 & \textbf{0.9149} \\
Both (Inf + Isch) & 124 & 0.9688 & 1.0000 & \textbf{0.9841} \\
\midrule
\textbf{Macro Average} & \textbf{1,191} & \textbf{0.9134} & \textbf{0.9279} & \textbf{0.9201} \\
\bottomrule
\end{tabular}
\end{table}

\subsection{Wagner Depth Ordinal Grading Results}
Table~\ref{tab:wagner_results} summarizes ordinal depth grading performance across the deduplicated Wagner validation cohort ($n = 303$). DFU-WAGNER net achieves an exact classification accuracy of \textbf{68.65\%} (208/303) and an adjacent-grade accuracy of \textbf{91.42\%} (277/303), bounding Mean Absolute Error to \textbf{0.4290} and attaining a Quadratic Weighted Kappa of \textbf{0.7765}.

\begin{table}[h]
\centering
\caption{Wagner Depth Ordinal Grading Performance (Validation Split, $n=303$)}
\label{tab:wagner_results}
\begin{tabular}{lcccc}
\toprule
\textbf{Evaluation Cohort} & \textbf{Accuracy (\%)} & \textbf{Adj. Acc. (\%)} & \textbf{MAE} & \textbf{QWK} \\
\midrule
Validation Split ($n=303$) & 68.65 & 91.42 & 0.4290 & 0.7765 \\
\bottomrule
\end{tabular}
\end{table}

\subsection{Ablation Study Results}
Table~\ref{tab:ablation} presents the systematic evaluation of all eight ablated configurations (A0 Baselines through A7) trained under controlled 25-epoch screening protocols. Each model's reported metrics correspond precisely to its saved best checkpoint under its designated monitoring metric.

\begin{table*}[t]
\centering
\caption{Systematic Ablation Study of DFU-WAGNER net Architectural and Regularization Components (Controlled 25-Epoch Screening Protocol)}
\label{tab:ablation}
\begin{tabular}{clccccc@{\hspace{1.5em}}cccc}
\toprule
\textbf{Exp ID} & \textbf{Configuration Description} & $\lambda_{\text{path}}$ & $\lambda_{\text{coral}}$ & $\lambda_{\text{lmcot}}$ & $\lambda_{\text{mmd}}$ & $\lambda_{\text{orth}}$ & \textbf{Macro-F1} & \textbf{Ischaemia F1} & \textbf{Wagner QWK} & \textbf{Wagner MAE} \\
\midrule
A0-Path & Pathology Single-Task Baseline & 1.0 (CBF) & 0.0 & 0.0 & 0.0 & 0.0 & 0.9306 & 0.9451 & --- & --- \\
A0-Wag & Wagner Single-Task Baseline & 0.0 & 2.5 & 0.0 & 0.0 & 0.0 & --- & --- & 0.7798 & 0.4554 \\
\textbf{A1} & \textbf{Full Proposed Multi-Task Net} & \textbf{1.0 (CBF)} & \textbf{2.5} & \textbf{0.5} & \textbf{0.1} & \textbf{0.1} & \textbf{0.9109} & \textbf{0.9348} & \textbf{0.7253} & \textbf{0.4818} \\
A2 & Without LMCoT Metric Learning & 1.0 (CBF) & 2.5 & 0.0 & 0.1 & 0.1 & 0.8995 & 0.9247 & 0.7268 & 0.4851 \\
A3 & Without UG-MMD Domain Alignment & 1.0 (CBF) & 2.5 & 0.5 & 0.0 & 0.1 & 0.8889 & 0.8866 & 0.7779 & 0.4554 \\
A4 & Without Orthogonality ($\mathcal{L}_{\text{orth}} = 0$)$^*$ & 1.0 (CBF) & 2.5 & 0.5 & 0.1 & 0.0 & 0.9121 & 0.9149 & 0.7387 & 0.4884 \\
A6 & Standard CE instead of CB-Focal & 1.0 (CE) & 2.5 & 0.5 & 0.1 & 0.1 & 0.9319 & 0.9556 & 0.7638 & 0.4653 \\
A7 & Uniform LMCoT (No Hard Mining) & 1.0 (CBF) & 2.5 & 0.5 & 0.1 & 0.1 & 0.8678 & 0.8400 & 0.7359 & 0.4917 \\
\bottomrule
\multicolumn{11}{l}{\footnotesize $^*$Note: In the experimental plan, the planned A5 ablation (``No Task Disentanglement'') was operationalized via A4 ($\lambda_{\text{orth}}=0$) to maintain tensor dimensional compatibility.}
\end{tabular}
\end{table*}

\subsection{Verified Computational Efficiency}
The exact architectural parameter counts and memory footprints---verified directly from the model implementation---are reported in Table~\ref{tab:param_breakdown} and Table~\ref{tab:compute_efficiency}:
\begin{itemize}
    \item \textbf{Total Architecture Parameters}: \textbf{59,783,374} (100\% trainable).
    \begin{itemize}
        \item CNN Backbone (EfficientNetV2-S): $20,177,488$ parameters (33.75\%).
        \item ViT Backbone (Swin-Transformer Tiny): $27,519,354$ parameters (46.03\%).
        \item Symmetrical Projections: $2,629,120$ parameters (4.40\%).
        \item TaskDisentangler: $8,139,520$ parameters (13.62\%).
        \item Task Heads (Classifier + Wagner CORAL): $1,317,892$ parameters (2.20\%).
    \end{itemize}
    \item \textbf{Memory Footprint}: Full FP32 weights require \textbf{228.06 MB}; FP16 weights require \textbf{114.03 MB}.
    \item \textbf{Peak Training VRAM}: Approximately \textbf{4,310 MB} under AMP FP16 at batch size 8, operating comfortably within standard 6 GB commercial GPU hardware constraints.
\end{itemize}

\begin{table}[h]
\centering
\caption{Detailed Layer-by-Layer Architectural Specifications and Parameter Counts}
\label{tab:param_breakdown}
\begin{tabular}{lccc}
\toprule
\textbf{Submodule} & \textbf{Output Dim} & \textbf{Parameter Count} & \textbf{\% Total} \\
\midrule
EfficientNetV2-S & 1,280 & 20,177,488 & 33.75\% \\
Swin-Transformer Tiny & 768 & 27,519,354 & 46.03\% \\
CNN Projection Head & 1,280 & 1,642,240 & 2.75\% \\
ViT Projection Head & 1,280 & 986,880 & 1.65\% \\
TaskDisentangler ($z_s, z_p, z_d$) & 768 / 512 / 512 & 8,139,520 & 13.62\% \\
Pathology Classifier Head & 4 & 658,948 & 1.10\% \\
Wagner Ordinal Head (CORAL) & 4 & 658,944 & 1.10\% \\
\midrule
\textbf{Total Architecture} & \textbf{Dual Tasks} & \textbf{59,783,374} & \textbf{100.00\%} \\
\bottomrule
\end{tabular}
\end{table}

\begin{table}[h]
\centering
\caption{Computational Efficiency and Hardware Footprint on NVIDIA RTX 3050 Laptop GPU}
\label{tab:compute_efficiency}
\begin{tabular}{lccc}
\toprule
\textbf{Configuration} & \textbf{Total Params} & \textbf{FP32 Size} & \textbf{Peak VRAM} \\
\midrule
Lightweight (B0 + MobileViT-S) & 12,541,200 & 47.84 MB & $\approx 1,450\text{ MB}$ \\
Optimized (B0 + Swin-Tiny) & 35,158,798 & 134.12 MB & $\approx 2,750\text{ MB}$ \\
\textbf{Primary Candidate (DFU-WAGNER net)} & \textbf{59,783,374} & \textbf{228.06 MB} & \textbf{$\approx 4,310\text{ MB}$} \\
\bottomrule
\end{tabular}
\end{table}

% =========================================================================
% SECTION VII: DISCUSSION
% =========================================================================
\section{Discussion}

\subsection{Minority-Class Ischaemia Representation}
In the DFUC2021 challenge benchmark \cite{cassidy2021diabetic}, detecting arterial ischaemia is a primary diagnostic hurdle because ischaemic wounds represent only 3.81\% of samples. Under standard cross-entropy, frequent categories dominate gradient updates, resulting in reported test F1-scores below 0.4200 \cite{galdran2021convolutional, toofanee2023dfu}. Clinically, unrecognized ischaemic ulceration delays vascular intervention, increasing the risk of major tissue loss \cite{armstrong1998validation}.

DFU-WAGNER net addresses this imbalance through Class-Balanced Focal Loss combined with Hard-Pair Tangent LMCoT Metric Learning on $z_{\text{shared}}$. On the validation cohort, the architecture achieves an \textit{Ischaemia} precision of \textbf{0.8958}, recall of \textbf{0.9348}, and F1-score of \textbf{0.9149}, identifying 43 out of 46 true ischaemic presentations. Concurrently, the compound \textit{Both} category achieves 100\% recall (124/124) with an F1-score of \textbf{0.9841}.

\subsection{Multi-Task Learning and Pareto Trade-Offs}
A key empirical insight from the ablation study is the existence of authentic multi-task Pareto trade-offs. In the 25-epoch screening run, the single-task pathology baseline (A0-Path) achieved a Macro-F1 of \textbf{0.9306}, whereas the full multi-task network (A1) achieved \textbf{0.9109}. Similarly, standard Cross-Entropy (A6) achieved a Macro-F1 of \textbf{0.9319} and a Wagner QWK of \textbf{0.7638}. 

This demonstrates that joint optimization across disparate clinical domains introduces gradient competition between biological surface erythema (pathology) and deep anatomical necrosis (depth). However, while single-task A0-Path predicts only pathology, the multi-task network concurrently predicts ordinal tissue invasion with 89.11\% adjacent accuracy. Extending training to 45 epochs with refined schedules narrows this gap, elevating multi-task Macro-F1 to \textbf{0.9201} while preserving dual-task clinical utility.

\subsection{Analysis of Architectural and Regularization Components}
\begin{enumerate}
    \item \textbf{Impact of Hard-Pair Mining (A1 vs. A7)}: Replacing active hardest-pair mining with uniform LMCoT sampling drops the DFUC Macro-F1 from \textbf{0.9109} to \textbf{0.8678} ($-4.31\%$) and reduces minority \textit{Ischaemia} F1 from \textbf{0.9348} to \textbf{0.8400} ($-9.48\%$). This confirms that active batch-level mining along clinical decision boundaries is essential for separating subtle ischaemic discoloration from uninfected granulation tissue.
    \item \textbf{Impact of Domain Alignment (A1 vs. A3)}: Disabling UG-MMD ($\lambda_{\text{mmd}} = 0.0$) degrades pathology Macro-F1 from 0.9109 to 0.8889 and ischaemia F1 to 0.8866. Interestingly, Wagner QWK increases to 0.7779 under A3, indicating that without alignment, the network favors the larger Wagner cohort at the expense of DFUC pathology. UG-MMD acts as a regularizing anchor ensuring balanced cross-cohort representation.
    \item \textbf{Impact of Metric Learning (A1 vs. A2)}: Disabling tangent LMCoT ($\lambda_{\text{lmcot}} = 0.0$) degrades Macro-F1 to 0.8995 and lowers overall accuracy to 86.40\%, demonstrating that hyperspherical metric clustering enforces regularizing intra-class compactness on the shared latent manifold.
    \item \textbf{Subspace Orthogonality (A1 vs. A4)}: Removing the orthogonality constraint ($\lambda_{\text{orth}} = 0$) results in slight DFUC Macro-F1 change (0.9121) while lowering Ischaemia F1 to 0.9149 and shifting Wagner QWK, confirming that unconstrained feature sharing permits subtle task interference.
\end{enumerate}

\subsection{Ordinal Biophysical Geometry Preservation}
By reformulating Wagner grading through Consistent Rank Logits (CORAL) \cite{cao2020rank}, DFU-WAGNER net respects the monotonic biophysical reality of ulcer invasion. In nominal multi-class formulations \cite{rostami2024deep, zhou2025deep}, a model that predicts Grade 0 for a Grade 4 gangrenous wound incurs the same nominal penalty as predicting Grade 1 for a Grade 0 wound. Under CORAL, four cumulative binary cutoffs enforce monotonic rank ordering. Consequently, our primary model bounds Mean Absolute Error to \textbf{0.4290} and achieves an adjacent-grade accuracy of \textbf{91.42\%} (with a Quadratic Weighted Kappa of \textbf{0.7765}), ensuring that diagnostic discrepancies remain localized to contiguous stages.

% =========================================================================
% SECTION VIII: LIMITATIONS
% =========================================================================
\section{Limitations}
To preserve scientific and clinical integrity, we declare the following limitations:
\begin{enumerate}
    \item \textbf{Absence of Patient Metadata}: Neither DFUC2021 nor the Wagner dataset provides patient identification tags, clinical case numbers, or longitudinal timestamps. While our cryptographic hash deduplication guarantees exact image-level independence, true inter-patient independence across splits cannot be independently verified.
    \item \textbf{Web-Origin Scraping Overlap}: Programmatic hash audit revealed 728 identical image pairs across public DFUC2021 and Wagner releases. While we have provided complete audit manifests and a leak-free split quarantine protocol for future training, researchers evaluating multi-task combinations of these datasets must account for this cross-cohort duplication.
    \item \textbf{Single-Seed Initialization}: Reported ablation runs were conducted under fixed random seed 42 to establish parameter-for-parameter comparability. Formal multi-seed variance quantification across randomized initializations (e.g., Seeds 123, 2026) remains for extended multi-center evaluations.
    \item \textbf{Evaluation Split vs. Hidden Test Set}: The reported 0.9201 Macro-F1 reflects a 20\% development validation split ($n = 1,191$), whereas official DFUC2021 challenge leaderboards reflect an independent hidden test set ($n \approx 5,734$). These numbers are not directly interchangeable.
    \item \textbf{Non-Causal Nature of Explainability}: Grad-CAM visual attribution heatmaps reflect post-hoc gradient sensitivity and must not be interpreted as histological tissue boundaries, segmentation masks, or definitive causal clinical evidence.
    \item \textbf{Clinical Deployment Status}: This work presents a computational research framework and is \textbf{not cleared for autonomous clinical diagnosis or unmonitored patient triage without direct expert podiatric supervision}.
\end{enumerate}

% =========================================================================
% SECTION IX: CONCLUSION
% =========================================================================
\section{Conclusion}
This paper presented \textbf{DFU-WAGNER net}, a multi-task deep representation learning framework designed to address clinical taxonomy fragmentation in automated diabetic foot ulcer assessment. By coupling an \textbf{EfficientNetV2-S} CNN branch with a \textbf{Swin-Transformer Tiny} branch through symmetrical projections, DFU-WAGNER net captures localized epidermal textures alongside macroscopic spatial foot morphology within a unified 2,560-dimensional representation. 

The three-way \textbf{TaskDisentangler} module partitions features into shared anatomical, pathology-private, and depth-private unit hyperspheres, regularized by a cosine \textbf{Orthogonality Loss} ($\mathcal{L}_{\text{orth}}$). Cross-cohort distribution shifts are minimized via \textbf{Uncertainty-Gated MMD} on the shared subspace, while \textbf{Class-Balanced Focal Loss}, \textbf{Hard-Pair Tangent LMCoT Metric Learning}, and \textbf{CORAL Ordinal Regression} optimize task-specific objectives. 

Evaluated across DFUC2021 and deduplicated Wagner cohorts, DFU-WAGNER net achieves a Macro-F1 of \textbf{0.9201}, overall accuracy of \textbf{90.18\%}, and minority \textit{Ischaemia} F1-score of \textbf{0.9149} (93.48\% sensitivity) on the development validation cohort, alongside an ordinal Wagner QWK of \textbf{0.7765} and adjacent depth accuracy of \textbf{91.42\%}. Systematic ablations confirm that active hard-pair mining and uncertainty-gated domain alignment prevent negative cross-cohort transfer and minority collapse. In addition, an exhaustive cross-dataset cryptographic audit provides the community with open manifests to resolve web-origin duplicate overlap. DFU-WAGNER net establishes an architecturally rigorous, computationally efficient foundation for multi-faceted diabetic wound characterization.

% =========================================================================
% REFERENCES
% =========================================================================
\bibliographystyle{IEEEtran}
\bibliography{references}

\end{document}
